\chapter{Slopes over the Boundary of Weight Space}
\label{chap:lwx-slopes}

We can now state what \cite{LWX} prove.  Their target is the Coleman--Mazur--Buzzard--Kilford
conjecture \cite[Conjecture 1.2]{LWX}, transposed from the eigencurve to the spectral curve
$\mathrm{Spc}_D$ of overconvergent automorphic forms for a definite quaternion algebra.  That
conjecture has three clauses: sufficiently close to the boundary,
\begin{enumerate}
  \item $\mathrm{Spc}^{>r}_D$ is a disjoint union of countably many connected components $Z_n$,
    each finite and flat over the annulus $\mathcal{W}^{>r}$ via the weight map;
  \item there are rationals $\alpha_1 \leq \alpha_2 \leq \cdots \to \infty$ with
    $\lvert a_p(z)\rvert = \lvert T_{\mathrm{wt}(z)}\rvert^{\alpha_n}$ for $z \in Z_n$;
  \item the sequence $\alpha_1, \alpha_2, \dots$ is a disjoint union of \emph{finitely many
    arithmetic progressions}, counted with multiplicity (at least for large indices).
\end{enumerate}
\cite[Theorem 1.3]{LWX} proves a coarse form of (1)--(2) over the full annulus
$\mathcal{W}^{>1/p}$: a decomposition indexed by the intervals $[n,n]$ and $(n,n+1)$, with
$v(a_p(x)) \in \varphi(q)\,v(T_{\mathrm{wt}(x)})\cdot I$ on the piece $X_I$, and explicit degrees.
\cite[Theorem 1.5]{LWX} proves (1)--(2) on the nose over a smaller annulus
$\mathcal{W}^{>\lambda}$, $\lambda = p^{-8/((p^{2}-1)t+8)}$, in the form
\[
  v\bigl(a_p(y)\bigr) \;=\; \varphi(q)\,v\bigl(T_{\mathrm{wt}(y)}\bigr)\,\alpha_i(\omega),
  \qquad y \in Y_{i,\omega},
  \tag{1.5.1}
\]
and then, in a second half beginning ``Moreover'', proves (3).  This is the \emph{spectral halo}:
approaching the boundary of weight space, the spectral curve resolves into pieces of a completely
explicit shape.

Theorem~\ref{thm:lwx318} gave a \emph{lower} bound polygon.  The content of this chapter is that,
under an explicit and checkable hypothesis, the bound is attained on a band around each of the
points $n_k = kpt$, and that this pins the whole polygon.  Three ingredients:
a combinatorial upper bound (\S1), a vertex analysis converting touching into $T$-free
arithmetic (\S2), and the analytic Claim of \cite[\S4.2]{LWX} which controls the indices in
between (\S3).  Section~4 assembles the polygon-level form of Theorem 1.5, and
\S\ref{sec:what-15-says} says precisely how much of \cite[Theorem 1.5]{LWX} that is.
Section~5 transports everything to the genuine Hecke operator, and
Section~\ref{sec:stepone} discharges the hypothesis (Steps I and III of \cite[\S3.23]{LWX}),
granted adelic Atkin--Lehner data, and proves the degree formulas of \cite[Theorem 1.3]{LWX},
\cite[Corollary 1.4]{LWX} and the second half of \cite[Theorem 1.5]{LWX} at the level of Newton
polygons.

\section{The upper bound polygon}

Everything here is $\N$-combinatorics of the exponent
$\lambda(n) = \sum_{k<n}\bigl(\lfloor k/t\rfloor - \lfloor k/pt\rfloor\bigr)$ of
Definition~\ref{def:lwxlambda}.

\begin{definition}
  \label{def:touchx}
  \uses{def:lwxlambda}
  \lean{LWX.touchX, LWX.two_mul_lwxLambda_touchX, LWX.lwxUpperTwice}
  \leanok
  Set $n_k = k\,p\,t$, the abscissa of the $k$-th \emph{touching vertex}.  Then
  \cite[(3.23.1)]{LWX}
  \[
    \lambda(n_k) \;=\; \frac{k^{2}p(p-1)t}{2},
  \]
  and the \emph{upper bound polygon} is the polygon through the points $(n_k, \lambda(n_k))$,
  linear of slope $(k+\tfrac12)(p-1)$ on $[n_k, n_{k+1}]$.  We work with twice its exponent, so
  as to stay in $\N$.
\end{definition}

\begin{lemma}[{\cite[Lemma 4.1]{LWX}}]
  \label{lem:lwx41}
  \uses{def:touchx}
  \lean{LWX.two_mul_lwxLambda_le_lwxUpperTwice,
        LWX.lwxUpperTwice_sub_two_mul_lwxLambda_le}
  \leanok
  The upper and lower polygons agree at every $n_k$ and differ by at most
  $(p^{2}-1)t/8$ in between:
  \[
    0 \;\leq\; \mathrm{upper}(n) - \lambda(n) \;\leq\; \frac{(p^{2}-1)t}{8}
    \qquad\text{for all } n .
  \]
\end{lemma}

So the two polygons are uniformly close; the whole strategy is to show the Newton polygon is
squeezed between them and touches the lower one at the $n_k$.

\begin{lemma}[The band identities, {\cite[p.~25]{LWX}}]
  \label{lem:band}
  \uses{def:touchx}
  \lean{LWX.lwxLambda_touchX_add, LWX.lwxLambda_touchX_sub, LWX.lwxLambda_touchX_add_ge,
        LWX.lwxLambda_touchX_sub_ge}
  \leanok
  $\lambda(n_k\pm i) = \lambda(n_k) \pm k(p-1)i$ for $0\leq i\leq t$, and beyond the band of
  width $t$ the excess is at least $\lvert i\rvert - t$.
\end{lemma}

The band identity is the reason the lower bound polygon is \emph{linear} of slope $k(p-1)$ across
$[n_k-t, n_k+t]$: the increments $\lfloor n/t\rfloor - \lfloor n/pt\rfloor$ are constant there.

\section{Vertices and Theorem 1.3}

Write $c_n = \sum_m b_{n,m}T^{m}$ as in \S\ref{chap:lwx-halo}, and recall from
Theorem~\ref{thm:lwx-sharp} that the polygon touches the lower bound at $n$ exactly when
$b_{n,\lambda(n)}$ is a unit.  This condition does not mention $T$.

\begin{definition}
  \label{def:unitband}
  \uses{thm:lwx-sharp, def:touchx}
  \lean{LWX.IsUnitCoeff, LWX.leftIndex, LWX.rightIndex, LWX.HasUnitBand}
  \leanok
  Say $n$ is a \emph{unit index} if $b_{n,\lambda(n)}\in\Zp^{\times}$.  Let
  $n_k^{-}$ be the least unit index in $[n_k-t, n_k]$ and $n_k^{+}$ the greatest in
  $[n_k, n_k+t]$.  The \emph{touching hypothesis} $\mathrm{HasUnitBand}(k)$ asserts that both
  exist.
\end{definition}

\begin{lemma}
  \label{lem:unitband-from-touching}
  \uses{def:unitband, thm:lwx318, thm:lwx-sharp}
  \lean{LWX.hasUnitBand_of_height_eq, LWX.height_specCharSeries_touchX}
  \leanok
  $\mathrm{HasUnitBand}(k)$ holds if and only if the Newton polygon of $\Char(P)(T_0)$ touches
  the lower bound polygon at $(n_k, \lambda(n_k)v(T_0))$ for one --- equivalently every --- halo
  point $T_0$.
\end{lemma}

This equivalence is the pivot of the whole chapter.  Touching is a statement at one weight; the
unit band is a statement about the integral matrix and holds for all weights at once.  So a
single weight at which one can verify touching --- which is what \cite[\S3.23 Step I]{LWX}
supplies from Atkin--Lehner theory and classicality --- propagates to the entire annulus.

\begin{theorem}[{\cite[Theorem 1.3]{LWX}}, the vertex reading]
  \label{thm:lwx13-vertices}
  \uses{def:unitband, thm:lwx318, thm:lwx-sharp, lem:band, def:isnewtonpolygonof}
  \lean{LWX.bandLine, LWX.height_specCharSeries_eq_bandLine,
        LWX.bandLine_lt_height_specCharSeries_of_lt_leftIndex,
        LWX.bandLine_lt_height_specCharSeries_of_rightIndex_lt}
  \leanok
  Assume $\mathrm{HasUnitBand}(k)$.  Then at every halo point $T_0$ the points
  $\bigl(n_k^{-}, \lambda(n_k^{-})v(T_0)\bigr)$ and $\bigl(n_k^{+}, \lambda(n_k^{+})v(T_0)\bigr)$
  are consecutive vertices of the Newton polygon of $\Char(P)(T_0)$, joined by the segment of
  slope $k(p-1)v(T_0)$ through $(n_k,\lambda(n_k)v(T_0))$; strictly outside $[n_k^-, n_k^+]$ the
  polygon lies strictly above that line.
\end{theorem}

\begin{proof}
  \leanok
  \uses{thm:lwx-sharp, lem:band, thm:lwx318, def:isnewtonpolygonof}
  On the band, Lemma~\ref{lem:band} says the lower bound polygon is the straight line of slope
  $k(p-1)v(T_0)$ through $(n_k, \lambda(n_k)v(T_0))$, and Theorem~\ref{thm:lwx-sharp} says the
  valuation of $c_n(T_0)$ equals its lower bound exactly at the unit indices --- in particular at
  $n_k^{\pm}$.  So the polygon meets the line at those two abscissae, hence coincides with it in
  between by convexity.  Outside the band, the excess of Lemma~\ref{lem:band} plus the definitive
  margin of Theorem~\ref{thm:lwx-sharp} at non-unit indices puts every point strictly above the
  line; the greatest-convex-minorant characterisation of
  Definition~\ref{def:isnewtonpolygonof} then gives the strict inequality for the polygon.
\end{proof}

\begin{corollary}[{\cite[(3.23.3)--(3.23.4)]{LWX}}]
  \label{cor:lwx13-slopes}
  \uses{thm:lwx13-vertices}
  \lean{LWX.unitSlope_specCharSeries_eq_iff, LWX.unitSlope_specCharSeries_mem_Ioo}
  \leanok
  The $j$-th slope of the polygon equals $k(p-1)v(T_0)$ exactly for
  $n_k^{-}\leq j < n_k^{+}$, and lies strictly between $k(p-1)v(T_0)$ and $(k+1)(p-1)v(T_0)$ for
  $n_k^{+}\leq j < n_{k+1}^{-}$.  In particular the multiplicities
  $\deg X_k = n_k^{+}-n_k^{-}$ and $\deg X_{(k,k+1)} = n_{k+1}^{-}-n_k^{+}$ are independent of
  $T_0$.
\end{corollary}

The independence is what \cite[Theorem 1.3]{LWX} needs for its component decomposition: the
pieces of the eigencurve over the annulus have degrees that do not vary.

\section{The Claim of \S4.2}

Corollary~\ref{cor:lwx13-slopes} controls the slopes on the bands and says the remaining slopes
lie strictly between consecutive values.  To pin those too, \cite[\S4.2]{LWX} restrict to a
smaller annulus, $0 < v(T) < 8/\bigl((p^{2}-1)t+8\bigr)$, and prove:

\begin{theorem}[{\cite[\S4.2, Claim]{LWX}}]
  \label{thm:lwx-claim}
  \uses{lem:lwx41, thm:lwx-sharp, def:touchx}
  \lean{LWX.unitIndex, LWX.IsBelowUpper, LWX.exists_isUnit_and_coeffVal_eq,
        LWX.isBelowUpper_of_coeffVal_lt, LWX.coeffVal_specCharSeries_eq_unitIndex_mul,
        LWX.le_coeffVal_specCharSeries_of_not_isBelowUpper}
  \leanok
  Let $m(l)$ be the least $m$ with $b_{l,m}\in\Zp^{\times}$.  If $(l, v(c_l(T_0)))$ lies strictly
  below the upper bound polygon for one $T_0$ in the small annulus, then $m(l)\geq\lambda(l)$,
  the same holds for every $T$ in the annulus, and
  \[
    v\bigl(c_l(T)\bigr) \;=\; m(l)\cdot v(T) \qquad\text{for all such } T .
  \]
  If it does not, then $v(c_l(T))\geq \mathrm{upper}(l)\cdot v(T)$ for every $T$.
\end{theorem}

\begin{proof}
  \leanok
  \uses{lem:lwx41, thm:lwx-sharp}
  This is the analytic core (4.2.1)--(4.2.4) of \cite{LWX}.  Being strictly below the upper bound
  forces, by Lemma~\ref{lem:lwx41}, that the dominant term in
  $c_l(T_0) = \sum_m \psi(b_{l,m})T_0^{m}$ occurs at an index $m$ with
  $\lambda(l)\leq m$ and $2m < \mathrm{upper}$ doubled; the smallness of the annulus makes the
  competition between the factors $p^{-v(b_{l,m})}$ and $\lVert T_0\rVert^{m}$ strict, so the
  dominant index is unique and is characterised $T$-freely as the least $m$ with $b_{l,m}$ a
  unit.  Uniqueness of the dominant term gives the valuation identity by the ultrametric
  principle, at every $T$ in the annulus simultaneously, since the characterisation of $m(l)$ no
  longer mentions $T$.
\end{proof}

Both the hypothesis and the conclusion have been made $T$-free: ``strictly below the upper bound
polygon'' becomes the predicate $2\,m(l) < \mathrm{upper}_{2}(l)$, both polygons scaling with
$v(T)$.  That is what allows the next section to define a single $T$-free polygon.

\section{Theorem 1.5}

\begin{definition}
  \label{def:shapepolygon}
  \uses{thm:lwx-claim, def:unitband, def:isnewtonpolygonof}
  \lean{LWX.shapeVal, LWX.shapePolygon, LWX.slopeRatio, LWX.isNewtonPolygonOf_shapePolygon}
  \leanok
  The \emph{shape polygon} is the Newton polygon of the $T$-free valuation sequence
  \[
    n \longmapsto
    \begin{cases}
      \lambda(n_k) & n = n_k, \\
      m(n) & n \text{ strictly below the upper polygon}, \\
      \infty & \text{otherwise},
    \end{cases}
  \]
  and $\mathrm{slopeRatio}(j)$ is its $j$-th unit slope.  It is well defined unconditionally.
\end{definition}

\begin{theorem}[{\cite[Theorem 1.5]{LWX}}, (1.5.1)]
  \label{thm:lwx15}
  \uses{def:shapepolygon, thm:lwx-claim, thm:lwx13-vertices, cor:ofslopes-transport}
  \lean{LWX.height_specCharSeries_eq_smul_shape,
        LWX.unitSlope_specCharSeries_eq_slopeRatio}
  \leanok
  Assume $\mathrm{HasUnitBand}(k)$ for every $k$.  Then for every $T$ in the small annulus
  $0 < v(T) < 8/\bigl((p^{2}-1)t+8\bigr)$ and every $j$, the $j$-th slope of the Newton polygon
  of $\sum_n c_n(T)X^{n}$, counted with multiplicity, is
  \[
    \mathrm{slope}_j(T) \;=\; v(T)\cdot\mathrm{slopeRatio}(j),
  \]
  with $\mathrm{slopeRatio}$ independent of $T$.  In \cite{LWX}'s normalisation
  $\mathrm{slopeRatio}(j) = \varphi(q)\,\tilde\alpha_j(\omega)$, the factor $\varphi(q)$ of
  (1.5.1) being absorbed into our ratio and $\tilde\alpha$ being the sequence $\alpha$ repeated
  according to the degrees $\deg Y_{i,\omega}$.
\end{theorem}

\begin{proof}
  \leanok
  \uses{thm:lwx-claim, thm:lwx13-vertices, def:shapepolygon}
  The Newton polygon of $\Char(P)(T)$ is the convex hull of the points $(n, v(c_n(T)))$.  At
  $n = n_k$ the height is $\lambda(n_k)v(T)$ by Theorem~\ref{thm:lwx13-vertices}; at $n$ strictly
  below the upper polygon it is $m(n)v(T)$ by Theorem~\ref{thm:lwx-claim}; and every other point
  lies above the upper polygon and so above the hull of those two families.  Thus the two hulls
  agree after scaling by $v(T)$, and reading unit slopes off a scaled polygon multiplies them by
  the same scalar.
\end{proof}

\begin{theorem}[The remaining clauses of {\cite[Theorem 1.5]{LWX}}]
  \label{thm:slope-growth}
  \uses{def:shapepolygon, thm:lwx318}
  \lean{LWX.monotone_slopeRatio, LWX.tendsto_slopeRatio_atTop}
  \leanok
  The ratios $\mathrm{slopeRatio}(0)\leq\mathrm{slopeRatio}(1)\leq\cdots$ increase, and tend to
  infinity.  Both statements are unconditional: they need neither the small annulus nor the
  touching hypothesis.
\end{theorem}

\begin{proof}
  \leanok
  \uses{thm:lwx318, def:shapepolygon}
  Monotonicity is convexity of the shape polygon.  For the growth: unit coefficients occur only
  at heights $\geq\lambda(n)$ by Theorem~\ref{thm:lwx318}, so the shape sequence dominates
  $\lambda$, whose increments $\lfloor k/t\rfloor-\lfloor k/pt\rfloor$ are unbounded.  Given $M$,
  the line $y = Mx - M^{2}pt$ lies on or below every point of the shape sequence, hence on or
  below the shape polygon by the supporting-line lemma; comparing with the height as the sum of
  the first $n$ slopes and using monotonicity gives a slope $\geq M-1$.
\end{proof}

The supporting-line lemma used here is a spec-level consequence of the greatest-convex-minorant
characterisation, recorded separately.

\begin{lemma}[Supporting lines]
  \label{lem:supporting-line}
  \uses{def:isnewtonpolygonof, def:ofslopes}
  \lean{IsNewtonPolygonOf.twoSlope_le_height, IsNewtonPolygonOf.line_le_height,
        NewtonPolygon₀.unitSlope_eq_of_height_eq, NewtonPolygon₀.monotone_toReal_unitSlope}
  \leanok
  A two-slope broken line lying on or below every point of an admissible valuation sequence lies
  on or below its Newton polygon; with equal slopes this is the supporting-line statement for a
  single line.  Moreover a unit slope is the corresponding height increment, and the unit slopes
  of a constructed polygon are monotone and never $\bot$.
\end{lemma}

\section{What Theorem 1.5 says, and what \S4 proves}
\label{sec:what-15-says}

Theorem~\ref{thm:lwx15} is a statement about Newton polygons; \cite[Theorem 1.5]{LWX} is a
statement about a rigid analytic space.  It is worth being exact about the gap.  \cite{LWX}'s
theorem reads:

\begin{quote}
  Let $\omega : \Delta \to \Zp^{\times}$ be a character.  Then there exists
  $\lambda \in (0,1)$ such that there exists a sequence of rational numbers
  $\alpha_0(\omega), \alpha_1(\omega), \dots$ in increasing order and tending to infinity such
  that $\mathrm{Spc}^{>\lambda}_{D,\omega}$ is a disjoint union
  $\coprod_{i\geq 0} Y_{i,\omega}$ of rigid analytic spaces finite and flat over
  $\mathcal{W}^{>\lambda}_\omega$ via $\mathrm{wt}$, such that
  \[
    v\bigl(a_p(y)\bigr) = \varphi(q)\,v\bigl(T_{\mathrm{wt}(y)}\bigr)\,\alpha_i(\omega)
    \tag{1.5.1}
  \]
  for every $y \in Y_{i,\omega}$.  More precisely, if the tame level is neat, then we can take
  $\lambda = p^{-8/((p^{2}-1)t+8)}$ for $p > 2$ \dots
  \emph{Moreover}, let $M$ be a positive integer so that $p^{-q/p^{M-1}(p-1)} > \lambda$
  \dots  Then if we extend the sequence $\alpha_0(\omega), \alpha_1(\omega),\dots$ into
  $\tilde\alpha_0(\omega), \tilde\alpha_1(\omega), \dots$ with each $\alpha_i(\omega)$ appearing
  with multiplicity $\deg Y_{i,\omega}$, then
  $\tilde\alpha_0(\omega), \tilde\alpha_1(\omega), \dots$ is a disjoint union of
  $(p-1)p^{M-1}t/2$ arithmetic progressions with (the same) common difference
  $\varphi(q)p^{M}/2q^{2}$.  More precisely,
  \[
    \tilde\alpha_{j + q^{-1}p^{M}t}(\omega\omega_0^{2})
      = \tilde\alpha_j(\omega) + \frac{p^{M}}{q^{2}}
    \qquad\text{for any } j \geq 0.
    \tag{1.5.2}
  \]
\end{quote}

So the theorem is really three assertions, and the arithmetic progressions are the third.

\paragraph{(a) The polygon identity.}  This is what Theorem~\ref{thm:lwx15} proves, and it is
\cite{LWX}'s own reduction: the proof in \cite[\S4.2]{LWX} opens by saying that, the tame level
being neat, ``it suffices to show that for $T \in \Cp$ with
$0 < v(T) < \tfrac{8}{(p^{2}-1)t+8}$, the ratios to $v(T)$ of the slopes (counted with
multiplicity) of the Newton polygon of $\sum_{n\geq0} c_n(T)X^{n}$ are independent of the choice
of $T$.''  Theorem~\ref{thm:lwx15} is that sentence, with the ratios named
$\mathrm{slopeRatio}(j) = \varphi(q)\tilde\alpha_j(\omega)$ and produced as the unit slopes of an
explicit $T$-free polygon (Definition~\ref{def:shapepolygon}), on exactly \cite{LWX}'s annulus.
Theorem~\ref{thm:slope-growth} adds the two clauses ``in increasing order and tending to
infinity''.  We prove it over any complete ultrametric field receiving $\Zp$ isometrically, not
only $\Cp$, and --- through Theorem~\ref{thm:lwx-slopes-seam} --- for the genuine $U_p$ rather
than for the integral matrix.  As stated it is conditional on the touching hypothesis
$\mathrm{HasUnitBand}$; Theorem~\ref{thm:lwx-unitband} discharges that hypothesis at every vertex
and on every disc, granted the adelic data of Definition~\ref{def:al-data}, so that for the genuine
$U_p$ (1.5.1) holds with no further hypothesis.

\paragraph{(b) The geometric decomposition.}  The passage from ``the slope ratios are
$T$-independent'' to ``$\mathrm{Spc}^{>\lambda}_{D,\omega} = \coprod_i Y_{i,\omega}$ with each
$Y_{i,\omega}$ finite flat over the annulus'' is not formalised, and is out of scope for this
thesis: the spectral curve is a rigid analytic space, and no rigid geometry is developed here.
What we do have is its fibrewise shadow.  Theorem~\ref{thm:lwx-spectral} identifies, at each
weight in the annulus, the zeros of $\Char(P)(T_0)$ with the reciprocal $U_p$-eigenvalues on
$S^{D,\dagger,m}_{\kappa_{T_0}}(U)$, so (1.5.1) holds pointwise for every point of the spectral
curve lying over the annulus; and Theorem~\ref{thm:halo-riesz} builds, at each halo vertex, the
finite locally free $U_p$-stable summand over the \emph{whole} annulus $\Lhalo[1/T]$ whose
Fredholm determinant is the corresponding polynomial factor --- which is the algebra underlying
the finite flat piece.  Assembling those into a rigid space, and identifying the summands with
connected components, is the missing step.

\paragraph{(c) The arithmetic progressions.}  This is the ``Moreover'' clause, (1.5.2).  It is
\emph{not} a consequence of Theorem~\ref{thm:lwx15}: (1.5.2) relates the sequence at the
nebentypus $\omega$ to the sequence at $\omega\omega_0^{2}$ --- two \emph{different} characters,
shifted by the square of the Teichm\"uller character --- whereas Theorem~\ref{thm:lwx15} is a
statement at a single fixed $\omega$, one weight disc at a time.  \cite[Remark 1.6(4)]{LWX} says
where it comes from: ``The second half of Theorem 1.5 follows from the first half by classicality
results and Atkin--Lehner theory'' (an argument found independently by Bergdall and Pollack).
That is the route formalised in Theorem~\ref{thm:lwx152}, at the level of Newton polygons.  The
Atkin--Lehner input of \cite[\S3.23 Step I]{LWX} is supplied at conductor $p^{h+1}$ for every
level $h \geq 1$ and at every classical weight at once, by \emph{families} of adelic data
(Definition~\ref{def:al-data}); the slopes of the classical subspace are identified without the
classicality theorem \cite[Proposition 2.15]{LWX}; and the slope reflection, (4.2.5),
(4.2.6) $=$ (1.5.2) and the arithmetic progressions are then bookkeeping in the characters of
$\Delta$.  The same comparison of two weight discs is behind \cite[Corollary 1.4]{LWX}, the
periodicity of $\deg X_{I,\omega}$ modulo $\varphi(q)/2$, which is formalised as
Corollary~\ref{cor:lwx14}.

\paragraph{Degrees.}  Corollary~\ref{cor:lwx13-slopes} gives the multiplicities as
$\deg X_k = n_k^{+} - n_k^{-}$ and $\deg X_{(k,k+1)} = n_{k+1}^{-} - n_k^{+}$, which is what the
vertex analysis produces and is enough for $T$-independence.  \cite[Theorem 1.3]{LWX} additionally
evaluates them, at a neat tame level, as
$\deg X_{n,\omega} = r_{\mathrm{ord}}(\omega^{-1}\omega_0^{2n-2}) + r_{\mathrm{ord}}(\omega\omega_0^{-2n})$
and $\deg X_{(n,n+1),\omega} = qt - r_{\mathrm{ord}}(\omega^{-1}\omega_0^{2n}) -
r_{\mathrm{ord}}(\omega\omega_0^{-2n})$ in terms of dimensions of ordinary subspaces.  Those
identifications are Step III, and they too involve several nebentypi at once.  They are formalised
at the coefficient level, with $r_{\mathrm{ord}}$ the largest index of a unit coefficient of the
characteristic series --- which is how \cite{LWX} prove their Theorem 3.19 --- as
Theorem~\ref{thm:lwx13-degrees}; only their reading as degrees of rigid analytic components is
not, as in (b).

\section{Reading it on the genuine $U_p$}

Chapter~\ref{chap:lwx-model} identified $\Char(P)(T_0)$ with $\det(1-X\,U_p)$ on
$S^{D,\dagger,m}$ at every analyticity level (Theorem~\ref{thm:lwx-prop217-h}).  Composing gives
each of the statements above for the genuine Hecke operator; nothing new is proved, and no
neatness hypothesis is needed, because none of the polygon statements uses one.

\begin{theorem}
  \label{thm:lwx-slopes-seam}
  \uses{thm:lwx-prop217-h, thm:lwx318, cor:lwx13-slopes, thm:lwx15,
        lem:unitband-from-touching}
  \lean{LWX.isBelow_newtonPolygon_discHeckeCharPowerSeries,
        LWX.hasUnitBand_of_height_discHeckeCharPowerSeries_eq,
        LWX.unitSlope_discHeckeCharPowerSeries_eq_iff,
        LWX.unitSlope_discHeckeCharPowerSeries_mem_Ioo,
        LWX.unitSlope_discHeckeCharPowerSeries_eq_slopeRatio,
        LWX.exists_level_unitSlope_discHeckeCharPowerSeries_eq_slopeRatio}
  \leanok
  For $\det(1-X\,U_p)$ on $S^{D,\dagger,m}$, at every point of the boundary annulus
  $p^{-1}<\lVert T_0\rVert<1$: the Newton polygon lies on or above the lower bound polygon of
  Theorem~\ref{thm:lwx318}; if it touches it at $n_k$ at one halo point then
  $\mathrm{HasUnitBand}(k)$ holds; the slope readings of Corollary~\ref{cor:lwx13-slopes} hold;
  and, on the small annulus, so does (1.5.1) of Theorem~\ref{thm:lwx15}.
\end{theorem}

The last item deserves emphasis: \emph{the touching hypothesis has to be verified only once, at a
single weight}, and Lemma~\ref{lem:unitband-from-touching} then supplies it everywhere.

\section{Theorem 1.3: Steps I and III}
\label{sec:stepone}

\cite[Theorem 1.3]{LWX} --- the statement that the spectral curve over the boundary annulus is a
disjoint union of pieces finite and flat over it, with the slope of $a_p$ on the piece $X_I$
lying in $\varphi(q)v(T)\cdot I$ and with explicit degrees --- needs the touching hypothesis to be
discharged.  \cite[\S3.23]{LWX} do this in three steps.  \emph{Step I} compares the slopes on the
whole space with those on the finite-dimensional classical subspace, and uses Atkin--Lehner
theory to pin their total; this is what supplies the touching.  \emph{Step II} is the vertex
analysis of \S2 above.  \emph{Step III} computes the degrees $\deg X_{I,\omega}$ in terms of
dimensions of ordinary subspaces, using Hida theory and the theta exact sequence.  All three steps
are formalised, at the level of characteristic series and Newton polygons, granted the adelic
data of Definition~\ref{def:al-data}; \S\ref{sec:ledger} records what each of them bought.

No step uses Jacquet--Langlands.  \cite{LWX}'s proofs of their Proposition 3.22 (Atkin--Lehner)
and Proposition 2.15 (classicality) both appeal to it; here the first is replaced by a
double-coset computation (Proposition~\ref{prop:al-identity}) and the second is not needed.  A
per-result audit of every use of Jacquet--Langlands, and of our avoidance of it, is maintained
alongside the sources.

\begin{definition}[The conjugate, partner and target nebentypi]
  \label{def:conjchar}
  \lean{LWX.invChar, LWX.teichChar, LWX.partnerChar, LWX.targetChar}
  \leanok
  Let $\omega$ be a character of $\Delta$, $\omega_0$ the Teichm\"uller character and $k \geq 0$,
  and let $(k,\psi)$ be a classical weight on the disc $\omega$, so that
  $\psi|_\Delta = \omega\omega_0^{-k}$.  The \emph{conjugate} nebentypus is
  $\omega^{-1} : u \mapsto \omega(u)^{-1}$; the \emph{partner} nebentypus
  $\omega^{-1}\omega_0^{2k}$ is the disc of the Atkin--Lehner partner $(k,\psi^{-1})$; and the
  \emph{target} nebentypus $\omega\omega_0^{-2k-2}$ is the disc of the theta target
  $(-k-2,\psi)$.  Since every statement of Chapters~\ref{chap:lwx-halo}--\ref{chap:lwx-slopes} is
  proved for an arbitrary $\omega$, these further instances cost nothing, and the same $U_p$-datum
  serves all of them, the nebentypus entering only through the entry streams.
\end{definition}

\begin{proposition}[Atkin--Lehner, reduced]
  \label{prop:atkin-lehner}
  \uses{lem:charpoly-pairing}
  \lean{LWX.atkinLehner, LWX.atkinLehnerConj, LWX.atkinLehner_mul_atkinLehnerConj,
        LWX.atkinLehnerConj_mem_Iw, LWX.roots_charpoly_atkinLehner,
        LWX.norm_roots_charpoly_atkinLehner}
  \leanok
  The element $w = \begin{pmatrix}0&1\\-p^{m}&0\end{pmatrix}$ has determinant $p^{m}$, normalises
  the Iwahori subgroup, inverts the nebentypus, and conjugates $U_p$ to $U'_p$ --- all expressed
  by the inverse-free identity $w\,\gamma^{w} = \gamma\,w$.  \emph{Granting} the operator
  identity $U_p\circ U'_p = p^{k+1}$ on the classical space,
  Lemma~\ref{lem:charpoly-pairing} gives \cite[Proposition 3.22]{LWX}: the $U_p$-slopes at
  nebentypus $\psi$ and at $\psi^{-1}$ pair up to $k+1$.
\end{proposition}

The proof of \cite[Proposition 3.22]{LWX} on \cite[p.~24]{LWX} applies Jacquet--Langlands and
classifies the local component as a principal series; neither is available in Lean.  The operator
identity is instead proved directly, from adelic data of the following kind, by expanding the
composite over the coset representatives of $U_p$.

\begin{definition}[Adelic Atkin--Lehner data]
  \label{def:al-data}
  \lean{LWX.AtkinLehnerData, LWX.AtkinLehnerFamily, LWX.AtkinLehnerDataH,
        LWX.AtkinLehnerFamilyH}
  \leanok
  \emph{Atkin--Lehner data} of conductor $p^{h+1}$ at a nebentypus $\psi$ consist of a section
  $\iota_p : \mathrm{GL}_2(\mathbb{Q}_p) \to G$ of the $p$-component such that the level contains
  the lifts of the Iwahori subgroup, the central element $p$ at $p$ is a global central element
  times an element of the level, and the Atkin--Lehner element of level $p^{h+1}$ normalises the
  disc-$0$ part of the level; together with a Hecke character restricting to $\psi$ on the level.
  A \emph{family} of such data is a single section with a Hecke character for every disc $\omega$
  and every weight $k$.  The $U_p$-representatives are built from the section alone, so one
  $U_p$-datum serves every classical weight of the family at once.
\end{definition}

\begin{proposition}[The Atkin--Lehner identity]
  \label{prop:al-identity}
  \uses{def:al-data, prop:atkin-lehner, prop:theta}
  \lean{LWX.ClassicalDiscForms, LWX.atkinLehnerEquiv, LWX.discHeckeCl_comp_atkinLehner,
        LWX.atkinLehnerHypothesis_of_atkinLehnerData, LWX.discHeckeClH_comp_atkinLehnerH,
        LWX.atkinLehnerHypothesis_of_atkinLehnerDataH}
  \leanok
  Grant Atkin--Lehner data of conductor $p^{h+1}$ at $\psi$.  On the classical forms
  $S^{D}_{k+2}(K^{p}\Iw_{p^{h+1}};\psi)$ --- in the disc model, the forms with values in the
  locally polynomial functions of degree $\leq k$ --- the Atkin--Lehner element induces an
  isomorphism $W$ onto the classical forms at $\psi^{-1}$, and
  \[
    U_p \circ W^{-1} \circ U_p^{(\psi^{-1})} \circ W \;=\; p^{k+1} .
  \]
  Hence the hypothesis of Proposition~\ref{prop:atkin-lehner} holds at the classical points, and
  the $U_p$-slopes at $\psi$ and at $\psi^{-1}$ pair up to $k+1$.
\end{proposition}

\begin{proof}
  \leanok
  \uses{prop:atkin-lehner}
  Expand the composite over the coset representatives of $U_p$.  The term with trivial
  representative is central and contributes $p^{k+1}$.  After conjugation by the Atkin--Lehner
  element every other term is a sum of the nebentypus over the cosets of $1 + p^{h+1}\Zp$ in
  $1 + p^{h}\Zp$, which vanishes because the conductor is exactly $p^{h+1}$.
  Proposition~\ref{prop:atkin-lehner} then gives the pairing.
\end{proof}

\begin{proposition}[The theta operator and the classical subspace]
  \label{prop:theta}
  \uses{def:discforms, def:blockmaps}
  \lean{LWX.thetaDisc, LWX.thetaBlock, LWX.thetaDisc_eq_zero_iff, LWX.thetaBlock_eq_zero_iff,
        LWX.finrank_locPolyDegSubmodule, LWX.bol, LWX.thetaDisc_comp_discHeckeBlock_of_autFactor,
        LWX.thetaBlock_comp_discHeckeBlockOp_of_autFactor}
  \leanok
  In the disc model an element \emph{is} the family of Taylor coefficients on each disc, so
  $\theta^{r}$ --- $r$-fold differentiation, \cite[\S7]{Bu04} --- is the shift
  $(a,j)\mapsto (j+1)\cdots(j+r)\,c(a,j+r)$: a diagonal of integers composed with a shift,
  bounded by $1$ and column-finite, hence continuous, and block-diagonal over the discs.  Its
  kernel at $r = k+1$ is exactly the locally polynomial functions of degree $\leq k$, which is
  the first sentence of \cite[Proposition 4]{Bu04}, of dimension $(k+1)p^{h}$ per disc --- the
  local half of \cite[(3.21.1)]{LWX}.  If the automorphy factors of two weights have the classical
  shape of exponent $k$ and the target shape of exponent $k+2$, with the same constants, then
  $\theta^{k+1}$ intertwines $U_p$ at the first weight with $p^{k+1}U_p$ at the second.
\end{proposition}

The intertwining is \cite[\S7]{Bu04}'s displayed identity: Bol's identity
$\partial^{r}\bigl(u^{j+r}L^{-(j+1)}\bigr) = (\det\gamma)^{r}\,(j+r)_{r}\,u^{j}L^{-(j+r+1)}$, for
$u$ and $L$ the numerator and the denominator of the M\"obius transformation of $\gamma$, summed
over the coset decomposition of $U_p$.  Its hypotheses are placed on the automorphy factors, not on
the characters: a version in which the two weights shared an arbitrary finite part is false, since
$\theta$ differentiates that part when it is not locally constant.  Nothing here depends on
Jacquet--Langlands: \cite[Proposition 4]{Bu04}'s first sentence and its small-slope half are
Jacquet--Langlands-free, and its converse half, which does use it, is never imported.

\begin{proposition}[The theta sequence, in determinant form]
  \label{prop:theta-exact}
  \uses{prop:theta, lem:diag-intertwine, def:conjchar}
  \lean{LWX.IsThetaExact, LWX.TargetData, LWX.charPowerSeries_comp_one_sub_truncation_eq_rescale,
        LWX.isThetaExact_classicalData, LWX.targetData_classicalPoint}
  \leanok
  At the classical point of weight $k$ on the disc $\omega$, the Fredholm determinant of $U_p$ on
  the complement of the classical subspace equals the Fredholm determinant of $p^{k+1}U_p$ at the
  theta target $(-k-2,\psi)$, a halo point on the disc of the target nebentypus
  $\omega\omega_0^{-2k-2}$.
\end{proposition}

This is all that \cite[\S3.23 Step III]{LWX} uses of the exact sequence
$0 \to S^{D}_{k+2} \to S^{D,\dagger}_{(k,\psi)} \xrightarrow{(d/dz)^{k+1}}
S^{D,\dagger}_{(-k-2,\psi)} \to 0$.  Writing $\theta^{k+1}$ as a diagonal composed with a shift,
and choosing a section of the shift, turns the intertwining of Proposition~\ref{prop:theta} into
a \emph{diagonal} intertwining on the complement, and Lemma~\ref{lem:diag-intertwine} gives the
identity of determinants.  The operator $\theta^{k+1}$ is not surjective on the Banach model ---
the inverse diagonal is unbounded --- so the right-exactness of the sequence is neither proved nor
needed.

\begin{proposition}[The finite factor]
  \label{prop:finite-factor}
  \uses{prop:det-summand, lem:charpolyrev}
  \lean{TateFredholm.charPowerSeries_eq_mul_polynomial, LWX.charPowerSeries_eq_mul_of_stable}
  \leanok
  If a compactoid $u$ commutes with an idempotent whose range is finite free, its Fredholm
  determinant splits as the reversed characteristic polynomial of the finite piece times the
  determinant of the complement.
\end{proposition}

Step I uses exactly this, the classical subspace being a finite-dimensional $U_p$-stable piece.

\begin{theorem}[{\cite[\S3.23 Step I]{LWX}}: the touching]
  \label{thm:lwx-stepone}
  \uses{prop:finite-factor, prop:atkin-lehner, prop:theta, thm:lwx318,
        lem:unitband-from-touching}
  \lean{LWX.ClassicalData, LWX.classicalData, LWX.isStepOneTouching_of_atkinLehnerHypothesis,
        LWX.hasUnitBand_of_atkinLehnerHypothesis}
  \leanok
  Let $T_0$ be the classical point of a weight $(k,\psi)$ with $\psi$ of conductor $p^{2}$, and
  suppose that the hypothesis of Proposition~\ref{prop:atkin-lehner} holds between $T_0$ and its
  Atkin--Lehner partner.  Then the Newton polygon of $\Char(P)(T_0)$ passes through
  $\bigl(n_{k+1}, \lambda(n_{k+1})v(T_0)\bigr)$; hence $\mathrm{HasUnitBand}(k+1)$.
\end{theorem}

\begin{proof}
  \leanok
  \uses{prop:finite-factor, prop:atkin-lehner, thm:lwx318, lem:unitband-from-touching}
  By Proposition~\ref{prop:finite-factor} the characteristic series at $T_0$ is the determinant of
  the complement times the reversed characteristic polynomial of $U_p$ on the classical subspace,
  of dimension $n_{k+1} = (k+1)pt$ by Proposition~\ref{prop:theta}; so the height of the polygon at
  $n_{k+1}$ is at most $v(\det U_p)$ on that subspace, and the same holds at the partner point.
  The hypothesis gives $v(\det U_p) + v(\det U'_p) = n_{k+1}(k+1)v(p) = 2\lambda(n_{k+1})v(T_0)$,
  the partner point having the same norm as $T_0$.  With the lower bound of
  Theorem~\ref{thm:lwx318} at both points, both inequalities are equalities, and
  Lemma~\ref{lem:unitband-from-touching} converts touching into $\mathrm{HasUnitBand}(k+1)$.
\end{proof}

\begin{theorem}[The touching hypothesis, discharged]
  \label{thm:lwx-unitband}
  \uses{thm:lwx-stepone, prop:al-identity, def:al-data, thm:lwx-slopes-seam}
  \lean{LWX.hasUnitBand_of_atkinLehnerData, LWX.hasUnitBand_of_atkinLehnerFamily,
        LWX.unitSlope_discHeckeCharPowerSeries_eq_slopeRatio_of_atkinLehnerFamily}
  \leanok
  Grant a family of Atkin--Lehner data of conductor $p^{2}$.  Then $\mathrm{HasUnitBand}(k)$ holds
  for every $k$ and on every disc $\omega$.  Consequently Corollary~\ref{cor:lwx13-slopes} and
  Theorems~\ref{thm:lwx15} and \ref{thm:lwx-slopes-seam} hold with the touching hypothesis replaced
  by the family of data; in particular (1.5.1) holds for the genuine $U_p$ at every analyticity
  level.
\end{theorem}

\begin{proof}
  \leanok
  \uses{thm:lwx-stepone, prop:al-identity}
  At $k = 0$ the band is automatic.  For $k+1$, the classical point of weight $k$ on the disc
  $\omega$ lies in the halo, and Proposition~\ref{prop:al-identity} at the member of the family
  of weight $k$ on the disc $\omega$ supplies the hypothesis of Theorem~\ref{thm:lwx-stepone}
  there; the $U_p$-datum does not depend on $(\omega,k)$ because the family has a single section.
\end{proof}

\begin{theorem}[{\cite[Theorem 1.3]{LWX}}, the degrees]
  \label{thm:lwx13-degrees}
  \uses{thm:lwx-unitband, prop:theta-exact, prop:al-identity, def:conjchar, cor:lwx13-slopes}
  \lean{LWX.ordDim, LWX.ordDim_le_card, LWX.degX, LWX.degXint, LWX.degX_zero, LWX.degX_succ,
        LWX.degXint_zero, LWX.degX_succ_classicalPoint, LWX.degX_succ_of_atkinLehnerData,
        LWX.degX_succ_of_atkinLehnerFamily, LWX.degXint_of_atkinLehnerFamily,
        LWX.degXint_pos_of_atkinLehnerFamily}
  \leanok
  Let $r_{\mathrm{ord}}(\omega)$ be the largest index $d$ such that the $d$-th coefficient of the
  characteristic series on the disc $\omega$ is a unit of the halo ring; it is at most $t$.  With
  $\deg X_{k,\omega} = n_k^{+} - n_k^{-}$ and $\deg X_{(k,k+1),\omega} = n_{k+1}^{-} - n_k^{+}$ as
  in Corollary~\ref{cor:lwx13-slopes}, $\deg X_{0,\omega} = r_{\mathrm{ord}}(\omega)$; and, granted
  a family of Atkin--Lehner data of conductor $p^{2}$, for every disc $\omega$ and every $n \geq 0$
  \[
    \deg X_{n+1,\omega} = r_{\mathrm{ord}}(\omega^{-1}\omega_0^{2n})
      + r_{\mathrm{ord}}(\omega\omega_0^{-2n-2}),
    \qquad
    \deg X_{(n,n+1),\omega} = pt - r_{\mathrm{ord}}(\omega^{-1}\omega_0^{2n})
      - r_{\mathrm{ord}}(\omega\omega_0^{-2n}) \;>\; 0 .
  \]
\end{theorem}

\begin{proof}
  \leanok
  \uses{thm:lwx-unitband, prop:theta-exact, prop:al-identity}
  At $n_{k+1}$ the unit band splits the degree into the left gap $n_{k+1} - n_{k+1}^{-}$ and the
  right gap $n_{k+1}^{+} - n_{k+1}$, and the finite factor splits the polygon: the classical
  slopes are at most $(k+1)v(p)$ and those of the complement at least $(k+1)v(p)$.  The left gap
  counts the classical slopes equal to $(k+1)v(p)$; the Atkin--Lehner pairing of
  Proposition~\ref{prop:al-identity} matches them with the classical slopes $0$ at the partner
  point, which are all its ordinary slopes, so the left gap is
  $r_{\mathrm{ord}}(\omega^{-1}\omega_0^{2k})$.  The right gap counts the slopes of the complement
  equal to $(k+1)v(p)$; by Proposition~\ref{prop:theta-exact} these are the ordinary slopes of the
  theta target, so the right gap is $r_{\mathrm{ord}}(\omega\omega_0^{-2k-2})$.  The interval
  degrees are $pt$ minus the two gaps at consecutive vertices, an exact subtraction because
  $n_k^{+} \leq n_{k+1}^{-}$; positivity is $r_{\mathrm{ord}} \leq t$ with $p \geq 3$.
\end{proof}

\begin{corollary}[{\cite[Corollary 1.4]{LWX}}]
  \label{cor:lwx14}
  \uses{thm:lwx13-degrees, def:conjchar}
  \lean{LWX.degX_succ_mul_teichChar_sq, LWX.degXint_mul_teichChar_sq,
        LWX.degX_succ_mul_teichChar_pow, LWX.degXint_mul_teichChar_pow, LWX.degX_succ_add_period,
        LWX.degXint_add_period}
  \leanok
  Granted a family of Atkin--Lehner data of conductor $p^{2}$,
  $\deg X_{I,\omega} = \deg X_{I+1,\omega\omega_0^{2}}$ for $I = (0,1), 1, (1,2), 2, \dots$, and
  hence $\deg X_{I,\omega}$ is periodic in $I$, on the same range, modulo
  $\varphi(q)/2 = (p-1)/2$.
\end{corollary}

\begin{proof}
  \leanok
  \uses{thm:lwx13-degrees}
  Substitute $(n,\omega) \mapsto (n+1,\omega\omega_0^{2})$ in the formulas of
  Theorem~\ref{thm:lwx13-degrees}: the partner and target nebentypi of
  Definition~\ref{def:conjchar} do not change.  Iterate, and use $\omega_0^{p-1} = 1$.
\end{proof}

The source's list of intervals starts at $(0,1)$ and omits $I = 0$, and rightly: the identity
fails there in general, since $\deg X_{0,\omega} = r_{\mathrm{ord}}(\omega)$ while
$\deg X_{1,\omega\omega_0^{2}} = r_{\mathrm{ord}}(\omega^{-1}\omega_0^{-2}) +
r_{\mathrm{ord}}(\omega)$.

\begin{theorem}[{\cite[Theorem 1.5]{LWX}}, second half, at the polygon level]
  \label{thm:lwx152}
  \uses{thm:lwx-unitband, prop:al-identity, def:al-data, def:conjchar, def:shapepolygon}
  \lean{LWX.atkinLehnerHypothesis_of_atkinLehnerDataH,
        LWX.slopeRatio_add_slopeRatio_partnerChar_of_atkinLehnerFamilyH,
        LWX.slopeRatio_partnerChar_succ, LWX.slopeRatio_mul_teichChar_sq,
        LWX.slopeRatio_mul_teichChar_pow, LWX.slopeRatio_add_period}
  \leanok
  Let $h \geq 1$ satisfy $(p^{2}-1)t + 8 < 8p^{h-1}(p-1)$, and grant families of Atkin--Lehner
  data of conductors $p^{2}$ and $p^{h+1}$ over a single section.  Then for every disc $\omega$,
  every $k \geq 0$, every $i < N = t(k+1)p^{h}$ and every $j \geq 0$,
  \begin{align*}
    \mathrm{slopeRatio}_{\omega^{-1}\omega_0^{2k}}(N-1-i) + \mathrm{slopeRatio}_{\omega}(i)
      &= (k+1)\,p^{h-1}(p-1), \\
    \mathrm{slopeRatio}_{\omega\omega_0^{2}}(j + p^{h}t)
      &= \mathrm{slopeRatio}_{\omega}(j) + p^{h-1}(p-1), \\
    \mathrm{slopeRatio}_{\omega}\bigl(j + \tfrac{p-1}{2}\,p^{h}t\bigr)
      &= \mathrm{slopeRatio}_{\omega}(j) + \tfrac{p-1}{2}\,p^{h-1}(p-1).
  \end{align*}
  With $M = h+1$ and $\mathrm{slopeRatio} = \varphi(q)\tilde\alpha$, the second line is (1.5.2),
  and the third says that $\tilde\alpha_0(\omega), \tilde\alpha_1(\omega), \dots$ is a disjoint
  union of $(p-1)p^{M-1}t/2$ arithmetic progressions with common difference
  $\varphi(q)p^{M}/2q^{2}$.
\end{theorem}

\begin{proof}
  \leanok
  \uses{thm:lwx-unitband, prop:al-identity}
  At a classical point of conductor $p^{h+1}$ in the slope-reading region, the classical subspace
  carries exactly the first $N$ slopes: they are at most $(k+1)v(p)$ by
  Proposition~\ref{prop:al-identity} and $\lVert U_p\rVert \leq 1$, while those of the complement
  are at least $(k+1)v(p)$, through the theta target.  The Atkin--Lehner pairing reverses them
  against the partner nebentypus, which with (1.5.1) at both discs
  (Theorem~\ref{thm:lwx-unitband}) is the first line.  The first line at $(\omega,k)$ and at
  $(\omega,k+1)$, subtracted, is \cite[(4.2.5)]{LWX}; taking the disc $\omega^{-1}\omega_0^{2k}$
  gives (4.2.6), the second line; and $\omega_0^{p-1} = 1$ iterates it into the third.
\end{proof}

\section{The dependency ledger}
\label{sec:ledger}

It is worth recording, once, which input bought which conclusion, since the statements have
genuinely different requirements and the obvious guess --- that Steps I and III together give
everything --- is not right.

\begin{center}
\small
\begin{tabular}{p{0.27\textwidth}p{0.40\textwidth}p{0.24\textwidth}}
  \hline
  \textbf{Target} & \textbf{Input required} & \textbf{Status here} \\
  \hline
  $\mathrm{HasUnitBand}(k)$ for all $k$; hence unconditional
    Corollary~\ref{cor:lwx13-slopes}, Theorems~\ref{thm:lwx15}, \ref{thm:slope-growth},
    \ref{thm:lwx-slopes-seam}
  & Step I at conductor $q^{2}$: Atkin--Lehner \cite[Prop.~3.22]{LWX}, the dimension count
    \cite[(3.21.1)]{LWX}, and the initial-segment comparison (classicality
    \cite[Prop.~2.15]{LWX} turns out not to be needed)
  & formalised, granted adelic data: Theorem~\ref{thm:lwx-unitband} \\
  \hline
  \cite[Thm.~1.3]{LWX}'s degree formulas, and \cite[Cor.~1.4]{LWX}
  & Step III: Hida's \cite[Cor.~3.21]{LWX} at the coefficient level, and the theta exact sequence
    $0 \to S^{D}_{k+2} \to S^{D,\dagger}_{(k,\psi)} \xrightarrow{(d/dz)^{k+1}}
     S^{D,\dagger}_{(-k-2,\psi)} \to 0$ in determinant form
  & formalised at the coefficient level, granted adelic data: Theorem~\ref{thm:lwx13-degrees},
    Corollary~\ref{cor:lwx14} \\
  \hline
  \cite[Thm.~1.5]{LWX}'s second half: (1.5.2) and the arithmetic progressions
  & Step I's Atkin--Lehner input again, but at conductor $p^{M}$ for every admissible $M$,
    together with statements ranging over the $\omega_0^{2}$-orbit of nebentypi
  & formalised at the polygon level, granted adelic data: Theorem~\ref{thm:lwx152} \\
  \hline
  The decomposition $\mathrm{Spc}^{>\lambda}_{D,\omega} = \coprod_i Y_{i,\omega}$, finite flat
  & rigid analytic geometry
  & out of scope for this thesis; see \S\ref{sec:what-15-says}(b) \\
  \hline
\end{tabular}
\end{center}

Two remarks on the third row, since it is the one that is easy to misjudge.  First, (1.5.2) is a
statement about $\tilde\alpha$, the slope-ratio sequence counted with multiplicity --- our
$\mathrm{slopeRatio}$ --- so it is a \emph{polygon-level} statement and needs none of the rigid
geometry of the fourth row.  Second, it does not follow from Steps I and III as such: it uses no
degree formula, but it needs the classical space and the Atkin--Lehner identity at
\emph{arbitrary} level $p^{M}$ rather than at $q^{2}$, and it compares the sequence at $\omega$
with the sequence at $\omega\omega_0^{2}$, whereas
Chapters~\ref{chap:lwx-halo}--\ref{chap:lwx-slopes} work at a single fixed $\omega$.

Concretely, five objects had to be built on top of those chapters, and all five now exist:

\begin{enumerate}
  \item the classical space $S^{D}_{k+2}(K^{p}\Iw_{p^{m}};\psi)$ itself, as the forms with values
    in the locally polynomial functions of degree $\leq k$, at conductor $p^{2}$ and at every
    conductor $p^{h+1}$, with its dimension (Proposition~\ref{prop:theta});
  \item the operator identity $U_p \circ W^{-1} U_p^{(\psi^{-1})} W = p^{k+1}$ on that space
    (Proposition~\ref{prop:al-identity}), which replaces the standing hypothesis of
    Proposition~\ref{prop:atkin-lehner};
  \item the theta operator's intertwining, in the form of Proposition~\ref{prop:theta}, and the
    theta sequence in determinant form (Proposition~\ref{prop:theta-exact});
  \item the initial-segment comparison, which needs no new theory: faces of Newton polygons add
    under products (Theorem~\ref{thm:np-product}), and the classical and complementary slopes are
    separated by $(k+1)v(p)$;
  \item a nebentypus-varying layer: the classical points of conductor $p^{h+1}$ and their halo
    weights, the partner and target nebentypi (Definition~\ref{def:conjchar}), families of adelic
    data (Definition~\ref{def:al-data}), and statements over the $\omega_0^{2}$-orbit
    (Corollary~\ref{cor:lwx14}, Theorem~\ref{thm:lwx152}).
\end{enumerate}

What is \emph{not} formalised is the adelic data itself: Definition~\ref{def:al-data} is carried as
a hypothesis.  For the finite ad\`eles of a definite quaternion algebra each of its fields is class
field theory or a local computation at $p$, but that instantiation has not been carried out.

A remark on the prime $2$.  \cite{LWX} \emph{state} Theorems~1.3 and~1.5 for $p = 2$ as well
(with $q = 4$, $\Delta = (\mathbb{Z}/4)^\times$, and the analyticity thresholds $m \ge 2$,
$M \ge 4$), but every proof in the paper is written for $p > 2$ only --- ``we assume $p > 2$, the
case $p = 2$ being similar'' (\S4.1 and \S4.2).  There is therefore no written argument to
transcribe for $p = 2$, and it is used essentially throughout this development from the
exponential and logarithm onwards; $p \ne 2$ is a standing hypothesis here, and the results
formalised are the cases the paper actually proves, which are narrower than its stated theorems.

Items (1)--(4) are Step I and Step III; item (5) is what the second half of
\cite[Theorem 1.5]{LWX} and \cite[Corollary 1.4]{LWX} need on top of them.

To summarise the status of \cite{LWX}: \S2 and \S3's estimates
(Theorems~\ref{thm:lwx316}, \ref{thm:lwx318}, \ref{thm:lwx-sharp}), the model identifications
(Theorems~\ref{thm:prop31}, \ref{thm:lwx-prop217}, \ref{thm:lwx-prop217-h}), the slope theorems
(Corollary~\ref{cor:lwx13-slopes}, Theorems~\ref{thm:lwx15}, \ref{thm:slope-growth},
\ref{thm:lwx-slopes-seam}), the touching hypothesis (Theorem~\ref{thm:lwx-unitband}), the degree
formulas of \cite[Theorem 1.3]{LWX} (Theorem~\ref{thm:lwx13-degrees}), \cite[Corollary 1.4]{LWX}
(Corollary~\ref{cor:lwx14}) and the second half of \cite[Theorem 1.5]{LWX}
(Theorem~\ref{thm:lwx152}) are formalised and free of \texttt{sorry}, for $p \neq 2$, granted the
adelic data of Definition~\ref{def:al-data}, and with no use of Jacquet--Langlands.  What remains
is the instantiation of that data for a definite quaternion algebra, and the passage from Newton
polygons to rigid analytic spaces.
